\section{Propagating the selected terminal birth to a moment bound}\label{ex-selected-birth-receipt}

This section records the finite-volume terminal receipt used on the public selected-history route. It starts from the calibrated family, its same-history multiscale data, and the birth polynomial made from a finite list of observables. The result is a mode-by-mode normalized moment estimate. It is a conditional receipt: it uses the selected-history service and migration data stated below, rather than deriving those upstream construction inputs.

Let $\mathcal V$ be a finite marker set and let $\mathsf B$ be a downward-closed family of allowed subsets containing the empty set and all singletons. Its elements are the census modes. For each $a\in\mathcal V$, choose a grade label $s_a$, a terminal site $o_N(a)$, and a motion $(\pi_a,\epsilon_a)$ consisting of a signed coordinate transformation and an independent conjugation flag. Let $O_a$ be the corresponding local cylinder observable, and define its terminal translate by
\[
 T_a(U)=\sigma_{\epsilon_a}\!\left(O_a\bigl(\operatorname{pull}_{\pi_a,o_N(a)}U\bigr)\right),
\]
where $\sigma_0$ is the identity and $\sigma_1$ is complex conjugation. The pullback includes both the signed transformation and the translation to $o_N(a)$.
Assume the selected family makes each $T_a$ continuous on the compact terminal configuration space. Write $R_N$ for the selected terminal representative and set
\[
 \mathcal B_N(U)=R_N(U)\,\operatorname{Pass}\!\left(
    \sum_{a\in\mathcal V} e_{\{a\}}T_a(U)\right),
\]
where the marker algebra has one basis element for each allowed mode and square-free multiplication: $e_S e_{S'}=e_{S\cup S'}$ if the sets are disjoint and their union is allowed, and $0$ otherwise. The polynomial pass is the finite exponential in this nilpotent algebra. For every allowed mode $I\in\mathsf B$, its coefficient satisfies
\[
 (\mathcal B_N(U))_I=R_N(U)\prod_{a\in I}T_a(U).
\]
To verify the identity, expand the $r$th power in the finite exponential. A term contributing to an allowed mode $I$ must use precisely the $|I|$ distinct singleton markers in $I$: repeated markers vanish, and no term of another degree can produce $I$. Every ordering of those singletons contributes the same product. Downward closure ensures that each intermediate union in such an ordering is allowed, so the $|I|!$ orderings cancel the factorial in the exponential coefficient. The empty mode comes from the zeroth power and has coefficient $R_N$. The centered-copy application uses the full census $\mathsf B=2^{\mathcal V}$; the general mode estimate below does not require that choice.
For this selected terminal representative, $R_N$ is the bare Wilson density. Write
\[
 Z_{N,p}:=\int R_N(U)\,d\mu_N(U)>0,
 \qquad
 \omega_{N,p}(F):=Z_{N,p}^{-1}\int R_N(U)F(U)\,d\mu_N(U).
\]

\begin{theorem}[Selected-history birth moment bound]\label{thm:ex-selected-birth-moment}
Fix a selected family, volume $N$, period $p$, grade $d$, and reference scale $k$ with $C.\mathrm{ref}(d)\leq k<N$. Assume all labels $s_a$ have grade $d$. Suppose the family satisfies its calibration conditions and the same-history finite-volume receipt supplies the terminal birth, its backward Wilson transitions, the scale-$k$ service resources, and the positive vacuum treasury. Assume also the geometric inputs at $k$: positive weights $w_a$, site assignments $o_k(a)$, an occupancy bound, the migration identity for the site assignments, and the separated-origin condition for every nontrivial mode. Suppose further that
\[
 w_a\,\mathfrak m_a\leq \mathfrak f_k\leq1\quad(a\in\mathcal V),
\]
where $\mathfrak m_a=\Gamma_k\,\mathrm{paid}_a$ is the column-weighted coefficient size of label $s_a$ at scale $k$, $\Gamma_k\geq1$ is the selected column count, and $\mathfrak f_k$ is the scale's small price. The factor $\Gamma_k$ is needed by the local coefficient and activity response; the unweighted payment alone is insufficient for that response estimate. Require the service receipt to give measurable mode densities $H_I$, a common scalar $\sigma$, and real costs $A_k,D_k$ such that its representative is nonnegative, its support exponent is nonnegative, and its reserve is $f=1$, with
\[
 w_I|H_I(U)|\leq \exp(\sigma+A_k),\quad
 Z_{N,p}\geq\exp(\sigma-D_k),\quad
 A_k+D_k\leq |\Lambda_k|\mathcal T_k+|\mathcal V|(x+y+2f),
\]
where $w_I=\prod_{a\in I}w_a$, $x=3\mathfrak f_k$, and $y=\mathfrak f_k$. The same-history receipt identifies $H_I$ with the scale-$k$ coefficient of the backward birth. Then for every census mode $I$,
\[
 \left|\omega_{N,p}\!\left(\prod_{a\in I}T_a\right)\right|
 \prod_{a\in I}w_a
 \leq \exp\!\left( |\Lambda_k|\,\mathcal T_k+6|\mathcal V|\right),
 \label{ex-selected-birth-bound}
\]
where $\omega_{N,p}$ is the normalized finite-volume Wilson expectation, $\Lambda_k$ is the scale-$k$ lattice, and $\mathcal T_k$ is the finite-volume vacuum treasury cost. The conclusion is conditional on the listed service, geometry, and paid-coefficient hypotheses.
\end{theorem}

\begin{proof}
First, $\mathcal B_N$ is integrable. The selected terminal representative is continuous, each $T_a$ is continuous by assumption, and the finite marker polynomial is a finite sum of products of these functions. Hence every coefficient of $\mathcal B_N$ is continuous. The terminal configuration space is compact, so each coefficient is integrable against the normalized product-Haar prior.

For $k\leq i\leq N$, let $\mathcal B_i$ be the coefficient polynomial obtained by the selected backward transition. It satisfies
\[
 \mathcal B_i=\operatorname{Push}_i(\mathcal B_{i+1}),
 \qquad \mathcal B_N\text{ as above},
\]
where the fine-to-coarse kernel is normalized in its output variable:
\[
 \int k_i(U,V)\,d\mu_i(V)=1\qquad\text{for each fine configuration }U.
\]
The density push is therefore
\[
 (\operatorname{Push}_iF)(V)=\int k_i(U,V)F(U)\,d\mu_{i+1}(U),
\]
where $U$ is fine and $V$ is coarse. For an integrable coefficient $F$, slice integrability and Fubini give the precise conservation identity
\[
 \int (\operatorname{Push}_iF)(V)\,d\mu_i(V)
 =\int F(U)\left(\int k_i(U,V)\,d\mu_i(V)\right)d\mu_{i+1}(U)
 =\int F(U)\,d\mu_{i+1}(U).
\]
For this terminal birth, $R_N$ is the bare Wilson density, so $Z_{N,p}=\int R_N\,d\mu_N$. The selected transition density is jointly measurable, nonnegative and uniformly bounded on the product of the finite-volume configuration spaces. Multiplying an integrable coefficient by that density gives an integrable function on the product of the priors; the same bound also gives integrability on every fixed coarse slice. Fubini therefore applies at each step. Since coefficient extraction is a continuous linear map on the finite-dimensional marker algebra, it commutes with this finite integral. Inductively, for each mode $I$,
\[
 \int (\mathcal B_k(U))_I\,d\mu_k(U)
 =\int (\mathcal B_N(U))_I\,d\mu_N(U)
 =Z_{N,p}\,\omega_{N,p}\!\left(\prod_{a\in I}T_a\right),
 \label{ex-selected-birth-integral}
\]
where $Z_{N,p}>0$ is the finite-volume Wilson partition function. The second equality uses the terminal birth identity and the definition of normalized expectation. This establishes the exact finite-volume transfer of the moment; no continuum limit or interchange of an infinite sum is involved.

At scale $k$, the same-history receipt yields a served marked representation of $\mathcal B_k$. Its finite response estimate gives at marker $a$ a coefficient bound by three times the column-weighted coefficient size $\mathfrak m_a$ and a total activity bound by $\mathfrak m_a$. Multiplying by $w_a$ and using $w_a\mathfrak m_a\leq\mathfrak f_k$ gives the coefficient and activity bills $x=3\mathfrak f_k$ and $y=\mathfrak f_k$. The factor three is the local coefficient bound from the response estimate; the activity bill is its unmultiplied support charge. The separated-origin condition and occupancy estimate are what make this finite response estimate valid for a whole mode, and migration transports those facts through the backward scales. Since $\mathfrak f_k\leq1$, these satisfy $x+y+2\leq6$.

Here are the precise conditional price assumptions used at the scale-$k$ service step. The marked representation is lawful, its representative is pointwise nonnegative, its support exponent is nonnegative, and its served reserve is $f=1$. For each mode $I$, the local density estimate is
\[
 w_I|H_I(U)|\leq \exp(\sigma+A_k),
 \qquad w_I=\prod_{a\in I}w_a,
\]
with $x=3\mathfrak f_k$ and $y=\mathfrak f_k$. The vacuum denominator bound is
\[
 Z_{N,p}\geq\exp(\sigma-D_k),
\]
and the receipts give
\[
 A_k+D_k\leq |\Lambda_k|\mathcal T_k
            +|\mathcal V|(x+y+2f),\qquad f=1.
\]
The reserve contributes the term $2f$; the coefficient $+2$ is not paid merely by the marker fee. Integrating the pointwise inequality yields
\[
 w_I\left|\int H_I\,d\mu_k\right|
 \leq \exp(\sigma+A_k).
\]
Divide by the positive denominator lower bound. The scalar $\sigma$ cancels:
\[
 w_I\frac{|\int H_I\,d\mu_k|}{Z_{N,p}}
 \leq \exp(A_k+D_k)
 \leq \exp\bigl(|\Lambda_k|\mathcal T_k+6|\mathcal V|\bigr),
\]
using $x+y+2f=4\mathfrak f_k+2\leq6$. Equation \eqref{ex-selected-birth-integral} identifies this normalized coefficient integral with the desired Wilson moment and proves the result.
\end{proof}

The terminal product, continuity, and exact backward integration have been proved above. The service certificate, its local coefficient envelope, the origin-counting response, and the finite-volume vacuum lower bound remain explicit inputs to this receipt; they are part of input I-C02-funded-record.

\subsection{The local denominator on the retained history}

\begin{lemma}[Quiet coefficient and partition function]
\label{lem:ex-selected-birth-denominator}
Fix the selected history at $(N,p)$. For each $i\le N$ let $A_i$
be its unmarked representative, and put $A_i^{\rm ev}(U)=\operatorname{eval}(A_i)(U)$.
Assume the selected Wilson identities
$A_N^{\rm ev}=R_N$ and $A_i^{\rm ev}=\operatorname{Push}_i A_{i+1}^{\rm ev}$ for $i<N$,
where the push is the normalized transition in
Section~\ref{ex-selected-birth-receipt}. Its kernel is that of step $i$
at the coupling $\beta_{N,p,i+1}$, the same coupling used by the
transfer and by the birth chain. Then
\[
 \int A_i^{\rm ev}\,d\mu_i=Z_{N,p}\quad\text{for all }i\le N,
 \qquad
 (\mathcal B_i)_\varnothing=A_i^{\rm ev}\quad\text{for }k\le i\le N,
\]
the range on which the chain $\mathcal B_i$ is defined.
\end{lemma}
\begin{proof}
At $N$, the constant coefficient of the finite exponential is one,
so $(\mathcal B_N)_\varnothing=R_N=A_N^{\rm ev}$. If the assertion holds at
$i+1$, extracting the quiet coefficient commutes with the normalized
push by the same finite-dimensional integration argument used in
\eqref{ex-selected-birth-integral}. Thus
\[
 (\mathcal B_i)_\varnothing
 =\operatorname{Push}_i(\mathcal B_{i+1})_\varnothing
 =\operatorname{Push}_i A_{i+1}^{\rm ev}=A_i^{\rm ev}.
\]
For the integral identity, each $A_{i+1}^{\rm ev}$ is integrable. At
$i+1=N$ it is continuous on a compact space. Below $N$ it is the push of
an integrable function by the bounded normalized kernel, by the argument
used for \eqref{ex-selected-birth-integral}. Integrating and using the
normalized transition mass at every step gives the claimed equality
with the positive terminal partition function.
\end{proof}

\paragraph{Common width and the vacuum cutoff.}
Fix an admitted grade $d$ and radius $R>0$. Let $g_{d,R}\ge1$
bound the coefficient norm on the closed grade ball, which is
compact in the finite-dimensional grade space. Using the
calibrated reference readout and the adjacent coefficient cone of
Section~\ref{ex-centered-freight-producer}, the coefficient at level
$k$ satisfies $q_k(v)\le g_{d,R}D_k$ throughout the ball, with
$D_k=\prod_{i<k}\Gamma_{i+1}$. Indeed the reference coefficient is
$v$ and each forward bound multiplies by at most the next column
count. Extending the product back to zero only increases it. Here
$\Gamma_i\ge1$ are the columns of the fixed schedule. Writing $a_k^{\rm mesh}$
for the mesh spacing, take the single positive weight
\[
 \omega_k=\frac{(a_k^{\rm mesh})^d\zeta_{d,k}}
                 {\Gamma_k g_{d,R}D_k}.
\]
Multiplication by the local column-weighted price
$\Gamma_k(a_k^{\rm mesh})^{-d}q_k(v)$ gives at most
$\zeta_{d,k}$ for every label in the ball. The weight is independent
of the terminal volume, period, site and particular label. Moreover
$\zeta_{d,k}\le1$. This follows from the coin condition
$\zeta_{d,k}\Gamma_k\le1$ of the purchased schedule, item~(R5) of
Section~\ref{sec:actual-marked-portfolio-calibration}, together with
$\Gamma_k\ge1$. So the small price $\mathfrak f_k=\zeta_{d,k}$ satisfies
the hypothesis $\mathfrak f_k\le1$ of
Theorem~\ref{thm:ex-selected-birth-moment}.

On a served census at $k<N$, use the local prices
$x=3\zeta_{d,k}$, $y=\zeta_{d,k}$ and reserve one in
\eqref{ex-centered-steer}. In the small-ball bound of
Section~\ref{ex-centered-construction} choose the radius
$s_k=\widehat\beta_k^{-1/2}$. The paid scale has
$\widehat\beta_k>1$ and the selected threshold is $s_k$ times an
exponential of a square, so $0<s_k\le1$ and $s_k$ is at most the
threshold. Write $E$ for the local excess cost, $\kappa$ for the
unmarked activity allowance, and $\beta$ for the local coupling.
The numerator pays per site $E+\kappa$ and the denominator pays
$E+\beta\,\operatorname{baseBudget}(s_k)+4d_G\log(16/s_k)$,
with the same scalar in both. To compare their sum with the selected
treasury, use the selected representation bounds: the local excess
cost is at most the regular allowance plus the loop price at its
coupling, the activity allowance is bounded by the scheduled forced
allowance, and $0\le\beta\le2\widehat\beta_k$. The loop price is an
affine function of the coupling with nonnegative slope, and the
base budget is nonnegative. Replacing the coupling in both costs
by $2\widehat\beta_k$ therefore gives at most the scheduled treasury
at radius $s_k$. The denominator in this calculation is
$\int\operatorname{eval}(A)\,d\mu_k$, where $A$ is the representative of
the served representation at scale $k$. The served census input is
taken over the history's own representative, that is $A=A_k$; this is
part of the conditional served-census input below. Hence, by
Lemma~\ref{lem:ex-selected-birth-denominator}, the denominator is
$\int A_k^{\rm ev}\,d\mu_k=Z_{N,p}$.
Thus, whenever the shared retention and transfer give the served
census and its local response on the actual geometry, the bound
\eqref{ex-selected-birth-bound} holds with this same width throughout
the grade ball. This argument uses the conditional served response
only at the opened scale and makes no new choice of history.
